\chapter{Probability II}\label{ch:iv:probability-ii}

``You roll a die until two sixes in a row appear. How many rolls on average?'' The candidate who writes two
states on the whiteboard (no six yet; one six) and one equation for each has the answer, 42, in two minutes;
the one who starts summing over sequences is still at it when the interviewer moves on. The second
probability chapter is about expectations and processes, and nearly every question in it yields to one of four
ideas: linearity with indicators, equations on states, martingales and stopping, and a little geometry of the
uniform and exponential laws.

\section{Expectation by indicators}

\begin{definition}[Indicator decomposition]\label{def:iv:probability-ii:indicators}
The \emph{indicator decomposition}\index{indicator decomposition} of a count $N$ writes it as a sum of
indicators, $N = \sum_i \mathbf 1_{A_i}$, so that $\E[N] = \sum_i \P(A_i)$ by linearity, whether or not the
events $A_i$ are independent.
\end{definition}

The method's power is that dependence does not matter for the mean. It matters for the variance, which needs the
pairwise probabilities $\P(A_i \cap A_j)$.

\begin{example}[Fixed points, runs, records]\label{ex:iv:probability-ii:indicators}
A random permutation of $n$ items has on average one fixed point: item $i$ is in place with probability $1/n$,
and there are $n$ items. A sequence of $n$ fair coin tosses has on average $1 + (n-1)/2$ runs: a new run starts
at toss $i \ge 2$ with probability $\tfrac12$. A sequence of $n$ independent draws from a continuous law has on
average $H_n = 1 + \tfrac12 + \dots + \tfrac1n$ records, since draw $i$ is the largest so far with probability
$1/i$; for ten draws, about 2.93.
\end{example}

\begin{example}[Distinct values]\label{ex:iv:probability-ii:distinct}
$n$ draws from $N$ equally likely values show on average $N\big(1 - (1 - 1/N)^n\big)$ distinct values:
value $v$ is missed by all draws with probability $(1-1/N)^n$. The same indicator gives the expected number of
empty buckets in a hash table.
\end{example}

\section{Recursions on states}

A waiting time for a pattern, a first passage in a chain, or the value of a game with a few positions is
found by first-step analysis (One Quant Book~4, chapter~8): condition on the first step, write one linear
equation per state, and solve.

\begin{method}[Setting up first-step equations]\label{met:iv:probability-ii:firststep}
\begin{enumerate}
\item Choose states that remember exactly what the future needs: for a pattern, the longest suffix of what has
been seen that is a prefix of the pattern.
\item For each state, write $e_s = 1 + \sum_t P_{st}\, e_t$ (expected steps) or $h_s = \sum_t P_{st}\, h_t$ (a
probability of absorption), with the absorbing states' values fixed.
\item Solve, and check a limiting case.
\end{enumerate}
\end{method}

\begin{example}[Two sixes in a row]\label{ex:iv:probability-ii:sixes}
With $e_0$ the expected further rolls with no six pending and $e_1$ with one: $e_0 = 1 + \tfrac16 e_1 + \tfrac56
e_0$ and $e_1 = 1 + \tfrac56 e_0$, so $e_0 = 42$ (\cref{fig:iv:probability-ii:states}). For coins the same method
gives 6 tosses for HH and 4 for HT: patterns that overlap themselves take longer to appear.
\end{example}

\begin{omfigure}
\begin{tikzpicture}[font=\footnotesize,>={Stealth},
  st/.style={circle,draw=omDef,fill=omDef!6,minimum size=1.1cm,align=center},
  ab/.style={circle,draw=omThm,fill=omThm!8,minimum size=1.1cm,align=center}]
\node[st] (s0) at (0,0) {no six};
\node[st] (s1) at (4,0) {one six};
\node[ab] (s2) at (8,0) {done};
\draw[->,thick,gray] (s0) to[bend left=20] node[above,black] {$\tfrac16$} (s1);
\draw[->,thick,gray] (s1) to[bend left=20] node[below,black] {$\tfrac56$} (s0);
\draw[->,thick,gray] (s1) -- node[above,black] {$\tfrac16$} (s2);
\draw[->,thick,gray] (s0) to[out=205,in=155,looseness=6] node[left,black] {$\tfrac56$} (s0);
\end{tikzpicture}

\omcaption{The two-state chain for ``two sixes in a row''. Each arrow is one roll; the expected number of rolls
from ``no six'' to ``done'' is 42. Exact solution by sympy in \texttt{iv\_prob2.py}.}\label{fig:iv:probability-ii:states}
\end{omfigure}

Races between patterns are solved the same way, with two absorbing states. Penney's game (1969) is the classic
case: for three-toss coin patterns, some patterns beat others with probability two thirds or more, and the
relation is not transitive.

\section{Martingales and stopping}

A fair game is a martingale (One Quant Book~4, chapter~1), and the optional stopping theorem says that a bounded
stopping rule cannot change its expected value. Interview questions use it in two ways: to compute a probability
in one line, and to spot a false claim that a stopping rule creates an edge.

\begin{proposition}[Gambler's ruin]\label{prop:iv:probability-ii:ruin}
A walk moves $+1$ with probability $p$ and $-1$ with probability $q = 1 - p$, starting at 0. The chance of reaching
$+b$ before $-a$ is $a/(a+b)$ if $p = \tfrac12$, and $\big(1 - (q/p)^a\big)/\big(1 - (q/p)^{a+b}\big)$ otherwise.
For $p = \tfrac12$ the expected duration from $x$ between 0 and $n$ is $x(n - x)$.
\end{proposition}

\begin{proof}
For $p = \tfrac12$ the position is a martingale; for $p \neq \tfrac12$ the process $(q/p)^{S_t}$ is; optional
stopping at the exit time (bounded in expectation) gives the probabilities. For the duration, $S_t^2 - t$ is a
martingale when $p = \tfrac12$.
\end{proof}

\begin{proposition}[Wald's identity]\label{prop:iv:probability-ii:wald}
If $X_1, X_2, \dots$ are i.i.d.\ with finite mean and $N$ is a stopping time with $\E[N] < \infty$ (or a count
independent of the $X_i$), then $\E\big[\sum_{i \le N} X_i\big] = \E[N]\,\E[X_1]$. If $N$ is independent of the
$X_i$, the variance of the sum is $\E[N]\Var(X_1) + \Var(N)\E[X_1]^2$.
\end{proposition}

A small edge compounds strongly in a ruin problem: with $p = 0.51$ and symmetric targets of ten units, the
chance of winning is about 0.60, not 0.51 (\cref{iq:iv:probability-ii:7}).

\section{Continuous problems and coupling}

Two facts about uniform and exponential variables answer most continuous questions. $n$ independent uniform
points on $[0,1]$ cut it into $n+1$ spacings that are exchangeable, each with mean $1/(n+1)$, so the $k$-th
smallest has mean $k/(n+1)$. Independent exponential clocks of rates $\lambda_1, \dots, \lambda_k$ ring first
at an exponential time of rate $\sum \lambda_i$, and clock $i$ rings first with probability $\lambda_i/\sum_j
\lambda_j$, independently of when.

\begin{example}[The inspection paradox]\label{ex:iv:probability-ii:inspection}
Five trades arrive at uniform times in an hour. The average gap between consecutive events (the ends of the
hour included) is 10 minutes, but the gap that contains a moment chosen at random has expected length $2/7$ of
an hour, about 17 minutes: longer gaps are more likely to be the ones you land in. Wait-time statistics measured
from random moments overstate the typical gap.
\end{example}

\begin{definition}[Coupling argument]\label{def:iv:probability-ii:coupling}
A \emph{coupling argument}\index{coupling argument} compares two random quantities by constructing them on
one probability space, each with its correct law, so that an inequality between them holds outcome by outcome;
the inequality then holds for their probabilities or expectations.
\end{definition}

\begin{example}[A better coin wins more often]\label{ex:iv:probability-ii:coupling}
Draw $U_1, \dots, U_{10}$ uniform on $[0,1]$ and let toss $i$ be a success for coin $p$ if $U_i < p$. With
$p = 0.6$ and $p' = 0.5$ every success of the second coin is a success of the first, so ``at least seven
successes'' for the second implies it for the first, and $\P(\ge 7 \mid 0.6) \ge \P(\ge 7 \mid 0.5)$ without
computing either (they are 0.382 and 0.172). The same construction proves that the ruin probability of
\cref{prop:iv:probability-ii:ruin} falls as $p$ rises.
\end{example}

\section{Worked answers}

\begin{example}[Adjacent pairs by indicators]\label{ex:iv:probability-ii:adjacent}
\emph{``Five buy and five sell orders are arranged in a uniformly random sequence. How many adjacent pairs have the same
side, on average?''} There are nine adjacent pairs. For any one of them, the chance that both positions hold the same
side is $(5 \times 4 + 5 \times 4)/(10 \times 9) = 4/9$, so by linearity the expected count is $9 \times 4/9 = 4$. The
pairs are not independent (a long run creates several same-side pairs at once) and the argument does not need them to
be. \emph{Check:} with one buy and one sell the formula gives $1 \times 0 = 0$, and indeed the single pair always differs.
\end{example}

\begin{example}[A recursion with a pattern]\label{ex:iv:probability-ii:total}
\emph{``You roll a die and add the results until the total is at least 3. How many rolls, on average?''} Let $e_k$ be
the expected number of rolls needed to reach at least $k$ from zero, with $e_k = 0$ for $k \le 0$. The first roll moves
you by $j$ with probability $1/6$: $e_k = 1 + \tfrac16\sum_{j=1}^6 e_{k-j}$. Then $e_1 = 1$, $e_2 = 1 + \tfrac16 = 7/6$,
and $e_3 = 1 + \tfrac16(7/6 + 1) = 49/36 \approx 1.36$. A pattern appears, $e_k = (7/6)^{k-1}$, and the candidate
should say how far it goes: subtracting consecutive equations gives $e_k - e_{k-1} = (e_{k-1} - e_{k-7})/6$, and $e_{k-7}
= 0$ while $k \le 7$, so the pattern holds up to $e_7 = (7/6)^6 \approx 2.52$ and breaks at $k = 8$. For large $k$ the
recursion settles at $e_k \approx k/3.5 + 10/21$: one roll per 3.5 of total, the mean roll, plus a constant for the
overshoot. Saying where a pattern must break is worth as much as finding it.
\end{example}

\section{Question bank}

\begin{interviewq}[$\star$ \iqroles{trader, researcher} \iqfirm{market maker}]\label{iq:iv:probability-ii:1}
Twenty orders are put into a queue in a uniformly random order. On average, how many end up in their original
position?
\end{interviewq}

\begin{interviewq}[$\star$ \iqroles{trader} \iqfirm{market maker}]\label{iq:iv:probability-ii:2}
How many rolls of a die on average until the first six? Until the second six (not necessarily consecutive)?
\end{interviewq}

\begin{interviewq}[$\star$ \iqroles{researcher, developer} \iqfirm{any}]\label{iq:iv:probability-ii:3}
A fair coin is tossed ten times. What is the expected number of runs (maximal blocks of equal faces)?
\end{interviewq}

\begin{interviewq}[$\star$ \iqroles{trader, researcher} \iqfirm{any}]\label{iq:iv:probability-ii:4}
Two independent uniform numbers on $[0,1]$: what are the expected values of the smaller and the larger?
\end{interviewq}

\begin{interviewq}[$\star$ \iqroles{trader, developer} \iqfirm{proprietary firm}]\label{iq:iv:probability-ii:5}
Fills arrive at venue A as a Poisson process of rate 3 a minute and at venue B at rate 1 a minute,
independently. What is the chance the next fill comes from A, and how long do you wait for it on average?
\end{interviewq}

\begin{interviewq}[$\star\star$ \iqroles{trader} \iqfirm{market maker}]\label{iq:iv:probability-ii:6}
You roll a die until two consecutive rolls show the same face, whatever it is. How many rolls on average? Why is
it so different from 42?
\end{interviewq}

\begin{interviewq}[$\star\star$ \iqroles{trader, risk} \iqfirm{proprietary firm}]\label{iq:iv:probability-ii:7}
Each bet wins one unit with probability 0.51 and loses one unit otherwise. You stop when you are ten units up or
ten units down. What is the chance you stop up?
\end{interviewq}

\begin{interviewq}[$\star\star$ \iqroles{researcher, trader} \iqfirm{any}]\label{iq:iv:probability-ii:8}
A symmetric random walk starts at 3 and stops at 0 or 10. How many steps does it take on average?
\end{interviewq}

\begin{interviewq}[$\star\star$ \iqroles{risk, trader} \iqfirm{bank}]\label{iq:iv:probability-ii:9}
A desk makes a Poisson number of trades a day with mean 40, independent of their P\&L, each with mean 0.3 and
variance 4 (thousands). Give the day's expected P\&L and its standard deviation.
\end{interviewq}

\begin{interviewq}[$\star\star$ \iqroles{developer, researcher} \iqfirm{any}]\label{iq:iv:probability-ii:10}
Twenty keys are hashed uniformly into 365 buckets. How many distinct buckets are used on average?
\end{interviewq}

\begin{interviewq}[$\star\star\star$ \iqroles{trader, researcher} \iqfirm{market maker}]\label{iq:iv:probability-ii:11}
A fair coin is tossed until either HHT or HTH appears. Which is more likely to appear first, and with what
probability? Which has the shorter expected waiting time on its own? Explain why the two answers point in
different directions.
\end{interviewq}

\begin{interviewq}[$\star\star\star$ \iqroles{trader, researcher} \iqfirm{proprietary firm}]\label{iq:iv:probability-ii:12}
You will receive five offers, one at a time, each an independent uniform number on $[0,1]$. You must accept or
reject each on the spot and you accept exactly one. What rule maximises the expected value accepted, and what is
that value? What is the threshold for the first offer?
\end{interviewq}

\begin{interviewq}[$\star\star\star$ \iqroles{researcher, developer} \iqfirm{systematic fund}]\label{iq:iv:probability-ii:13}
Five orders arrive at independent uniform times in an hour. What is the expected arrival time of the second? A
monitor checks at a uniformly random moment and measures the gap between the order before and the order after
(or the hour's ends). What is the expected length of the gap it sees, and why does it differ from the average
gap?
\end{interviewq}

\begin{interviewq}[$\star\star\star$ \iqroles{researcher, risk} \iqfirm{systematic fund}]\label{iq:iv:probability-ii:14}
Prove without computing either probability that ten tosses of a coin with $p = 0.6$ give at least seven heads
more often than ten tosses of a fair coin. Then compute both.
\end{interviewq}

\begin{omsources}
\item W.~Penney, \emph{Journal of Recreational Mathematics}, October 1969, 241 (Penney's game).
\item A.~Wald, ``On cumulative sums of random variables'', \emph{Annals of Mathematical Statistics} 15(3),
1944, 283--296.
\item T.~Lindvall, \emph{Lectures on the Coupling Method}, Wiley, 1992.
\item One Quant Book~4, chapters 1, 6 and~8 (martingales; Poisson processes; Markov chains and first-step
analysis).
\end{omsources}
